% !PRINTED 298
\weberpage{298}
%<<CONT
Theiler $Q'$ von $Q$, der sowohl von $Q$ selbst als von der Einheitsgruppe verschieden ist. Dieser Theiler $Q'$ muss aber ein Normaltheiler von $P$ sein, weil, wenn $\pi$ in $Q$, $\gamma$ in $P$ enthalten ist, $\gamma^{-1}\pi\gamma$, was vom selben Grade wie $\pi$ ist, gleichfalls in $Q$ enthalten sein muss, und also zu $Q'$ geh\"ort, wenn $\pi$ zu $Q'$ geh\"ort. Da nun aber, wie oben bemerkt, $Q$ keinen echten Theiler haben kann, der gr\"osser als die Einheitsgruppe und zugleich Normaltheiler von $P$ ist, so muss $\lambda = 1$ sein.

\begin{center}
\S.\,76.\\
\gesperrt{Darstellung der Abel'schen Gruppe $Q$.}
\end{center}

Die Betrachtungen des vorigen Paragraphen haben gezeigt, dass in einer metacyklischen Gruppe $P$, die als \gesperrt{Galois}'sche Gruppe einer irreduciblen primitiven Gleichung $f(x) = 0$ des Grades $n = p^k$ auftreten kann, eine \gesperrt{Abel'sche Gruppe} $Q$ als Normaltheiler enthalten sein muss, deren Grad $p^k$ ist, und die ausser dem Einheitselemente nur Elemente vom Grade $p$ enth\"alt.

Nach \S.\,9 l\"asst sich diese Gruppe $Q$ durch eine Basis darstellen, die aus $k$ Elementen $A_1, A_2, \ldots, A_k$ vom Grade $p$ besteht, so dass jedes Element von $Q$ die Form erh\"alt:
\weq{1}{A_1^{z_1} A_2^{z_2} \ldots A_k^{z_k},}
und hierin durchlaufen $z_1, z_2, \ldots, z_k$ von einander unabh\"angig je ein volles Restsystem nach dem Modul $p$. Wir gehen nun, um die entsprechende Permutationsgruppe zu finden, von einer beliebigen Wurzel $x$ von $f(x)$ aus, bezeichnen die Wurzel, in die $x$ durch die Permutation (1) \"ubergeht, mit
\weq{2}{[z_1, z_2, \ldots, z_k],}
und setzen fest, dass dies Zeichen seine Bedeutung nicht \"andern soll, wenn die Zahlen $z_1, z_2, \ldots, z_k$ beliebig um Vielfache von $p$ ver\"andert werden. Der Wurzel $x$, von der wir ausgingen, kommt dann das Zeichen $[0, 0, \ldots, 0]$ zu, und durch das Zeichen (2) ist jede Wurzel von $f(x)$ ein und nur einmal dargestellt.

Wenn $x$ durch eine Permutation $A'$ in $x'$ und $x'$ durch $A''$ in $x''$ \"ubergeht, so geht $x$ durch die Permutation $A' A''$ in $x''$ \"uber. Daraus folgt aber, dass $[z_1, z_2, \ldots, z_k]$ durch die Permutation
\weq{3}{A = A_1^{\alpha_1} A_2^{\alpha_2} \ldots A_k^{\alpha_k}}
%CONT>>
